\documentclass[11pt,a4paper]{article}
\usepackage[utf8]{inputenc}
\usepackage{amsmath}
\usepackage{amssymb}
\usepackage{algorithm}
\usepackage{algpseudocode}  % 使用algpseudocode代替algorithmic
\usepackage{tikz}
\usetikzlibrary{shapes.geometric, arrows.meta, positioning}
\usepackage[margin=2cm]{geometry}

\title{MACPO: Multi-Agent Cooperative Parallel Optimization \\ with Reinforcement Learning}
\author{}
\date{}

\begin{document}

\maketitle

\section{Algorithm Flowchart}

\begin{center}
\begin{tikzpicture}[
    node distance=0.7cm,
    box/.style={rectangle, draw, minimum width=2cm, minimum height=0.6cm, font=\scriptsize},
    round/.style={rectangle, rounded corners, draw, fill=red!20, minimum width=2cm, minimum height=0.6cm, font=\scriptsize},
    proc/.style={box, fill=blue!15},
    deci/.style={diamond, draw, fill=green!15, aspect=2.5, minimum width=1.5cm, font=\scriptsize},
    comm/.style={box, fill=orange!15},
    out/.style={box, fill=yellow!15, minimum width=1.5cm, minimum height=0.5cm},
    arrow/.style={->, >=stealth, thick}
]

% Main flow
\node[round] (start) {Init: $\alpha=0$};
\node[deci, below=of start] (loop) {iter$<$max?};
\node[proc, below=of loop] (pso) {1:PSO/LLSO};
\node[proc, below=of pso] (best) {2:localBest};
\node[proc, below=of best] (conf) {3:conflict};
\node[deci, below=of conf] (gate) {4:CI$>$T?};

% Right branch
\node[comm, right=2cm of gate] (nego) {5:Exchange};
\node[comm, below=of nego] (nash) {6:Nash};
\node[comm, below=of nash] (grad) {7:Gradient};

% Left branch
\node[proc, left=1.5cm of gate] (skip) {8:Skip};

% Merge
\node[proc, below=1.2cm of gate] (merge) {9:globalBest};
\node[proc, below=of merge] (eval) {10:Evaluate};
\node[proc, below=of eval] (mpi) {11:MPI};

% Right output
\node[out, right=0.5cm of mpi] (fp) {$f_{pure}$};
\node[out, below=0.2cm of fp] (fpen) {$f_{penalty}$};

% Continue
\node[proc, below=0.8cm of mpi] (imp) {12:improve};
\node[proc, below=of imp] (rew) {13:reward};
\node[proc, below=of rew] (rl) {14:RL-$\alpha$};
\node[proc, below=of rl] (outp) {15:Output};
\node[deci, below=of outp] (check) {16:End?};
\node[round, below=of check] (end) {End};

% Arrows
\draw[arrow] (start) -- (loop);
\draw[arrow] (loop) -- node[right,font=\tiny]{Y} (pso);
\draw[arrow] (pso) -- (best);
\draw[arrow] (best) -- (conf);
\draw[arrow] (conf) -- (gate);
\draw[arrow] (gate) -- node[above,font=\tiny]{Y} (nego);
\draw[arrow] (nego) -- (nash);
\draw[arrow] (nash) -- (grad);
\draw[arrow] (grad) -| (merge);
\draw[arrow] (gate) -- node[above,font=\tiny]{N} (skip);
\draw[arrow] (skip) |- (merge);
\draw[arrow] (merge) -- (eval);
\draw[arrow] (eval) -- (mpi);
\draw[arrow, dashed] (mpi) -- (fp);
\draw[arrow] (fp) -- (fpen);
\draw[arrow] (mpi) -- (imp);
\draw[arrow] (imp) -- (rew);
\draw[arrow] (rew) -- (rl);
\draw[arrow] (rl) -- (outp);
\draw[arrow] (outp) -- (check);
\draw[arrow] (check) -- node[right,font=\tiny]{Y} (end);
\draw[arrow] (check.west) -| ++(-2.5,0) |- node[near end,above,font=\tiny]{N} (loop.west);
\draw[arrow] (loop.east) -| ++(1.5,0) |- node[near start,above,font=\tiny]{N} (end.east);

\end{tikzpicture}
\end{center}

\section{Algorithm Pseudocode}

\begin{algorithm}
\caption{MACPO with RL-based Adaptive Penalty}
\begin{algorithmic}[1]
\Require Global objective $F(x) = \sum_{i=1}^N f_i(x_i)$, Agent count $N$, Max evaluations $E_{\max}$
\Ensure Optimal solution $x^*$, Optimal value $F(x^*)$

\State \textbf{Initialize on each agent $i \in \{1, \ldots, N\}$:}
\State \quad Population $\mathcal{P}_i$ with $M=300$ particles, each $x_k \in \mathbb{R}^{D_i}$
\State \quad Consensus reference $x_{i,\text{con}} \gets \text{random solution}$
\State \quad Penalty coefficient $\alpha_i \gets 0$ \Comment{Start from zero, RL will adjust}
\State \quad RL policy network $\pi_\theta$: weights $W \in \mathbb{R}^{1 \times 2}$, bias $b \in \mathbb{R}$ (random init)
\State \quad Conflict history buffer $H_i \gets []$ (max size 10)
\State \quad Replay buffer $\mathcal{B}_i \gets \emptyset$
\State \quad Previous fitness $f_{\text{prev}} \gets \infty$

\While{$\text{eval\_count} < E_{\max}$}
    \State \textbf{/* PHASE 1: Local Optimization with Penalty */}
    \For{$g = 1$ to $\text{gen\_times}$}
        \For{each particle $x_k \in \mathcal{P}_i$}
            \State Compute conflict: $c_k \gets \sum_{d \in O_i} \frac{|x_k^d - x_{i,\text{con}}^d|}{x_{\max} - x_{\min}}$
            \State Evaluate penalized fitness: $h_i(x_k) \gets f_i(x_k) + \alpha_i \cdot c_k$
        \EndFor
        \State Update velocities and positions (standard PSO/LLSO)
    \EndFor
    
    \State \textbf{/* PHASE 2: Select Local Best */}
    \State Find best particle: $x_{i,\text{best}} \gets \arg\min_{x_k \in \mathcal{P}_i} h_i(x_k)$
    
    \State \textbf{/* PHASE 3: Compute Conflict for Monitoring */}
    \State $\text{conflict}_i \gets \sum_{d \in O_i} \frac{|x_{i,\text{best}}^d - x_{i,\text{con}}^d|}{x_{\max} - x_{\min}}$
    \State Append to history: $H_i.\text{push}(\text{conflict}_i)$
    
    \State \textbf{/* PHASE 4: Communication Decision (Three-Layer Gating) */}
    \State Smooth conflict: $\text{CI}_{\text{smooth}} \gets 0.7 \cdot \text{CI}_{\text{smooth}} + 0.3 \cdot \text{conflict}_i$
    \State Check conditions:
    \State \quad $\text{cond}_1 \gets (\text{iter} - \text{last\_comm} \geq 10)$ \Comment{Min interval}
    \State \quad $\text{cond}_2 \gets (\text{CI}_{\text{smooth}} > \text{threshold})$ \Comment{Adaptive threshold}
    \State \quad $\text{cond}_3 \gets (\text{rand}() < f_{\text{phase}})$ \Comment{Phase-based sampling}
    \State $\text{should\_comm}_i \gets \text{cond}_1 \land \text{cond}_2 \land \text{cond}_3$
    \State Global decision: $\text{do\_comm} \gets \text{MPI\_Allreduce}(\text{should\_comm}_i, \text{MAX})$
    
    \If{$\text{do\_comm} = \text{true}$}
        \State \textbf{/* PHASE 5-7: Negotiation */}
        \State \textit{Step 5:} Exchange solutions with neighbors via MPI\_Isend/Recv
        \State \textit{Step 6:} Nash bargaining on overlapping variables
        \State \quad Initialize $x_{i,\text{temp}} \gets x_{i,\text{best}}$
        \For{each neighbor $j$ and shared dimension $d \in O_i \cap O_j$}
            \State Evaluate cross-impact: $f_{ii} \gets f_i(x_{i,\text{best}})$, $f_{ij} \gets f_j(x_{i,\text{best}})$
            \State \quad\quad\quad\quad\quad\quad\quad $f_{ji} \gets f_i(x_{j,\text{best}})$, $f_{jj} \gets f_j(x_{j,\text{best}})$
            \If{$f_{ii} + f_{ji} > f_{ij} + f_{jj}$}
                \State $x_{i,\text{temp}}^d \gets x_{j,\text{best}}^d$ \Comment{Adopt neighbor's value}
            \EndIf
        \EndFor
        \State \textit{Step 7:} Gradient-based conflict detection (optional, can be skipped)
        \State Update consensus: $x_{i,\text{con}} \gets x_{i,\text{temp}}$
        \State $\text{last\_comm} \gets \text{iter}$
    \Else
        \State \textbf{/* PHASE 8: Skip Negotiation */}
        \State $x_{i,\text{con}} \gets x_{i,\text{best}}$ \Comment{Use local best directly}
    \EndIf
    
    \State \textbf{/* PHASE 9: Update Optimizer State */}
    \State Set global best in optimizer to $x_{i,\text{con}}$
    \State Re-evaluate all particles in $\mathcal{P}_i$ with updated $x_{i,\text{con}}$
    \State Sort $\mathcal{P}_i$ by penalized fitness $h_i$
    
    \State \textbf{/* PHASE 10: Compute Global Metrics */}
    \State Local pure fitness: $f_i^{\text{local}} \gets f_i(x_{i,\text{con}})$
    \State Local penalized fitness: $h_i^{\text{local}} \gets h_i(x_{i,\text{con}})$
    \State Global pure fitness: $f^{\text{pure}} \gets \text{MPI\_Allreduce}(f_i^{\text{local}}, \text{SUM}) / N$
    \State Global penalized fitness: $h^{\text{penalty}} \gets \text{MPI\_Allreduce}(h_i^{\text{local}}, \text{SUM}) / N$
    \State Penalty term: $\text{penalty} \gets h^{\text{penalty}} - f^{\text{pure}}$
    
    \State \textbf{/* PHASE 11: RL Update for Penalty Coefficient */}
    \State Compute improvement: $\text{improv} \gets (f_{\text{prev}} - f^{\text{pure}}) / \max(|f_{\text{prev}}|, 10^{-8})$
    \State Compute reward (Scheme 5): $r \gets \tanh\!\left(-\Delta\text{conflict}_{\text{rate}} \cdot \Delta f_{\text{pure,rate}} \cdot 100\right)$
    \State Build state: $s \gets [\text{conflict}_i, H_i[\text{end}] - H_i[\text{end}-1]]^T$ \Comment{[conflict, trend]}
    \State Policy forward: $a \gets \tanh(W \cdot s + b)$
    \State Update $\alpha$: $\alpha_i \gets \text{clip}(\alpha_i + 0.01 \cdot a, 0.01, 1.0)$
    \State Store experience: $\mathcal{B}_i.\text{add}((s, a, r, |r|))$ \Comment{Priority = |r|}
    \If{$|\mathcal{B}_i| \geq 32$}
        \State Sample 32 experiences by priority
        \For{each sampled $(s', a', r', p)$}
            \State Target action: $a_{\text{target}} \gets a' + 0.1 \cdot r'$
            \State Compute loss: $\mathcal{L} \gets (a_{\text{target}} - \pi_\theta(s'))^2$
            \State Gradient descent: $\theta \gets \theta - 0.01 \cdot \nabla_\theta \mathcal{L}$
        \EndFor
    \EndIf
    \State Update: $f_{\text{prev}} \gets f^{\text{pure}}$
    
    \If{$\text{rank}_i = 0$}
        \State Output metrics: (iter, eval\_count, $f^{\text{pure}}$, penalty, $\alpha_i$, conflict, improv, $r$)
    \EndIf
\EndWhile

\State \Return $x^* \gets \bigcup_{i=1}^N x_{i,\text{con}}$, $F(x^*) = N \cdot f^{\text{pure}}$

\end{algorithmic}
\end{algorithm}

\section{Mathematical Formulas by Flowchart Steps}

\subsection{Step 0: Initialization}
\textbf{Global objective function} to minimize (distributed across $N$ agents):
\begin{equation}
\min_{x} F(x) = \sum_{i=1}^N f_i(x_i), \quad x_i \in \Omega_i \subset \mathbb{R}^{d_i}
\end{equation}

\textbf{Notation convention:}
\begin{itemize}
    \item $F(x)$: Global objective (sum of all agents' local objectives)
    \item $f_i(x_i)$: Agent $i$'s local objective (pure fitness, no penalty)
    \item $h_i(x_i)$: Agent $i$'s penalized fitness (used in PSO optimization)
    \item Initial penalty coefficient: $\alpha_0 = 0$ (starts from zero, RL adjusts dynamically)
\end{itemize}

\textbf{Summary of notation system:}
\begin{center}
\begin{tabular}{|l|l|l|}
\hline
\textbf{Symbol} & \textbf{Meaning} & \textbf{Usage} \\
\hline
$F(x)$ & Global objective & $F(x) = \sum_{i=1}^N f_i(x_i)$ \\
$f_i(x_i)$ & Pure local fitness (agent $i$) & No penalty, true objective \\
$h_i(x_i)$ & Penalized local fitness (agent $i$) & $h_i = f_i + \alpha \cdot \text{conflict}_i$ \\
$\text{conflict}_i(x_i)$ & Conflict penalty (L1-norm sum) & $\sum_d \frac{|x_{i,d} - x_{i,\text{con}}^d|}{\text{range}}$ \\
$x_{i,\text{con}}$ & Consensus reference (agent $i$) & Updated each iteration (Step 9) \\
$f^{\text{pure}}$ & Average pure fitness & $\frac{1}{N}\sum_{i=1}^N f_i(x_{i,\text{con}})$ \\
$h^{\text{penalty}}$ & Average penalized fitness & $\frac{1}{N}\sum_{i=1}^N h_i(x_{i,\text{con}})$ \\
\hline
\end{tabular}
\end{center}

\textbf{Key points about conflict calculation:}
\begin{itemize}
    \item \textbf{Not averaged:} $\text{conflict}_i = \sum_{d \in O_i} \frac{|Δx_d|}{\text{range}}$ (sum over dimensions)
    \item \textbf{Range:} $\text{conflict}_i \in [0, |O_i|]$ scales with number of overlapping dimensions
    \item \textbf{RL compensation:} The learned $\alpha$ automatically adjusts for problem-specific conflict scales
    \item \textbf{Typical magnitude:} $\text{conflict}_i \approx 0.5 \sim 5.0$ for well-tuned systems
\end{itemize}

\textbf{Key relationship:} At convergence, $h_i(x^*) = f_i(x^*)$ and $F(x^*) = N \cdot f^{\text{pure}}$

\subsection{Step 1: PSO/LLSO/CSO Local Optimization}
\textbf{Penalized local fitness} for each agent $i$ (used in PSO):
\begin{equation}
h_i(x_i) = f_i(x_i) + \alpha_i(t) \cdot \text{conflict}_i(x_i)
\end{equation}

where the conflict penalty term is defined as:
\begin{equation}
\text{conflict}_i(x_i) = \sum_{d \in O_i} \frac{|x_{i,d} - x_{i,\text{con}}^d|}{x_{\max} - x_{\min}}
\end{equation}

\textbf{Important note:} The conflict is the \textbf{sum} over all overlapping dimensions (not averaged). This design choice means:
\begin{itemize}
    \item Larger overlap size $|O_i|$ → larger conflict value (proportional scaling)
    \item The RL-adjusted $\alpha$ will adaptively compensate for different problem scales
    \item Typical conflict magnitude: $\text{conflict}_i \in [0, |O_i|]$ (not normalized to [0,1])
\end{itemize}

\textbf{Symbol definitions:}
\begin{itemize}
    \item $h_i(x_i)$: Penalized fitness for agent $i$ (PSO optimization objective)
    \item $f_i(x_i)$: Pure local fitness for agent $i$ (no penalty)
    \item $x_i \in \mathbb{R}^{D_i}$: Local solution vector for agent $i$ (including all variables)
    \item $x_{i,d}$: The value of the $d$-th decision variable in agent $i$'s current solution $x_i$
    \item $O_i$: Set of overlapping variable indices that agent $i$ shares with neighbors
    \item $x_{i,\text{con}}^d$: \textbf{Consensus reference value} for dimension $d$ (from previous iteration's Step 9)
    \item $\alpha_i(t) \in [0.01, 1.0]$: Time-varying penalty coefficient adjusted by RL
    \item $\text{conflict}_i(x_i)$: Normalized conflict measure (L1-norm sum, range $[0, |O_i|]$)
\end{itemize}

\textbf{Physical meaning:} The conflict penalty $\text{conflict}_i(x_i)$ measures the \textbf{sum of normalized L1-distances} across all overlapping dimensions, driving overlapping variables towards agreement. The conflict scales linearly with $|O_i|$ (number of overlapping dimensions), and the RL mechanism adaptively adjusts $\alpha$ to maintain appropriate penalty strength.

\textbf{Example:} Consider a 1000-dimensional problem with 2 agents:
\begin{itemize}
    \item Agent 1: manages variables $x_1, x_2, \ldots, x_{600}$ (dimensions 1-600)
    \item Agent 2: manages variables $x_{401}, x_{402}, \ldots, x_{1000}$ (dimensions 401-1000)
    \item Overlapping region: $O_1 = O_2 = \{401, 402, \ldots, 600\}$ (200 shared variables)
\end{itemize}
For agent 1, the penalized fitness is:
\begin{equation}
h_1(x_1) = f_1(x_1) + \alpha_1(t) \cdot \text{conflict}_1(x_1)
\end{equation}
where
\begin{equation}
\text{conflict}_1(x_1) = \sum_{d=401}^{600} \frac{|x_{1,d} - x_{1,\text{con}}^d|}{x_{\max} - x_{\min}}
\end{equation}

\textbf{Numerical example:} Assume $x_{\max} - x_{\min} = 100$ and average deviation is $|x_{1,d} - x_{1,\text{con}}^d| = 1.0$:
\begin{equation}
\text{conflict}_1 = \sum_{d=401}^{600} \frac{1.0}{100} = 200 \times 0.01 = 2.0
\end{equation}

\textbf{Scaling observation:} With 200 overlapping dimensions, the conflict magnitude is $\sim 2.0$ (not 0.01). The RL mechanism learns to adjust $\alpha$ accordingly to maintain effective penalty strength.

\textbf{Important clarification:} In Step 1, the penalty is applied to \textbf{every particle} during PSO optimization:
\begin{itemize}
    \item Population size: $M = 300$ particles
    \item Each particle $k$: $x_k \in \mathbb{R}^{600}$ for agent 1
    \item Penalized fitness for particle $k$: $h_1(x_k) = f_1(x_k) + \alpha_1(t) \cdot \text{conflict}_1(x_k)$
    \item PSO updates all 300 particles using $h_i$, not $f_i$
\end{itemize}

\subsection{Step 2: Get Local Best (localBestPar)}
Extract the best solution from local population after PSO optimization:
\begin{equation}
x_i^{\text{best}} = x_{\text{bestIndex}}, \quad \text{bestIndex} = \arg\min_{k \in [1,M]} h_i(x_k)
\end{equation}

\textbf{Implementation:} \texttt{Optimizer->getBestPar()} returns the position vector of the particle with lowest penalized fitness $h_i$:
\begin{itemize}
    \item For PSO: Returns \texttt{swarm[bestParIndex]->X} (best particle's position)
    \item For LLSO/CSO: Returns \texttt{swarm[0]->X} (first particle after sorting)
    \item Result: $x_i^{\text{best}} \in \mathbb{R}^{D_i}$ is a complete solution vector, NOT an average
\end{itemize}

\textbf{Key distinction:}
\begin{itemize}
    \item Step 1: Evaluates \textbf{all $M$ particles} using $h_i(x_k)$ (penalized fitness)
    \item Step 2: Selects the \textbf{single best particle} with minimum $h_i$
    \item The selection is based on $h_i$, but later we also compute pure fitness $f_i(x_i^{\text{best}})$
\end{itemize}

\subsection{Step 3: Conflict Calculation}
Compute normalized conflict measure using the best solution obtained from Step 2:
\begin{equation}
\text{conflict}_i = \sum_{d \in O_i} \frac{|x_{i,\text{best}}^d - x_{i,\text{con}}^d|}{x_{\max} - x_{\min}}
\end{equation}

\textbf{Key points:}
\begin{itemize}
    \item Uses $x_{i,\text{best}}$ from Step 2 (the single best particle, not the entire population)
    \item $x_{i,\text{best}}^d$ denotes the value at dimension $d$ in this best solution
    \item $x_{i,\text{con}}^d$ is the consensus reference from previous iteration (Step 9)
    \item Normalized by variable bounds (dividing by $x_{\max} - x_{\min}$)
    \item Uses L1-norm (absolute value) for interpretable distance measurement
    \item Result: $\text{conflict}_i \in [0, |O_i|]$ represents sum of normalized deviations across overlapping dimensions
    \item \textbf{No averaging:} The conflict is summed over all dimensions, not divided by $|O_i|$
\end{itemize}

\textbf{Design rationale: Why sum instead of average?}

The current implementation uses $\text{conflict}_i = \sum_{d \in O_i} \frac{|Δx_d|}{\text{range}}$ (sum) rather than $\frac{1}{|O_i|}\sum_{d \in O_i} \frac{|Δx_d|}{\text{range}}$ (average). This design choice has important implications:

\begin{enumerate}
    \item \textbf{Problem-scale awareness:}
    \begin{itemize}
        \item With sum: $\text{conflict}_i \propto |O_i|$ (scales with overlap size)
        \item With average: $\text{conflict}_i \in [0,1]$ (fixed range regardless of problem size)
        \item Benefit: The penalty naturally reflects the complexity of coordination (more overlapping dimensions → higher coordination challenge)
    \end{itemize}
    
    \item \textbf{Adaptive α compensation:}
    \begin{itemize}
        \item The RL mechanism learns an appropriate $\alpha$ that compensates for the conflict scale
        \item For large $|O_i|$: RL learns smaller $\alpha$ to avoid over-penalization
        \item For small $|O_i|$: RL learns larger $\alpha$ to maintain sufficient penalty strength
        \item Result: $\alpha \times \text{conflict}_i$ adapts to maintain consistent penalty-to-fitness ratio
    \end{itemize}
    
    \item \textbf{Comparison with original paper:}
    \begin{itemize}
        \item Original MACPO: $\alpha = |f| / 512$ (fixed formula)
        \item Current implementation: $\alpha = \text{base\_alpha} \times \text{RL\_ratio}$ where $\text{base\_alpha} = |f| / 512$
        \item Both approaches scale $\alpha$ with objective magnitude, ensuring penalty remains meaningful
        \item Original: conflict includes $|O_i|$ in denominator, $\alpha$ is larger
        \item Current: conflict excludes $|O_i|$ in denominator, $\alpha$ is smaller (RL-adjusted)
        \item \textbf{Key insight:} The product $\alpha \times \text{conflict}$ should have similar order of magnitude in both versions
    \end{itemize}
    
    \item \textbf{Numerical example:}
    \begin{itemize}
        \item Problem: $|O_i| = 200$, average deviation = 0.5, range = 100, $f = 10^5$
        \item Current (sum): $\text{conflict} = 200 \times (0.5/100) = 1.0$
        \item Original (average): $\text{conflict} = (200 \times 0.5/100) / 200 = 0.005$
        \item Current: $\alpha \approx (10^5 / 512) \times 0.001 = 0.195$ (RL-adjusted) $\Rightarrow$ penalty $\approx 0.195$
        \item Original: $\alpha = 10^5 / 512 = 195.3$ $\Rightarrow$ penalty $\approx 195.3 \times 0.005 = 0.977$
        \item Both achieve penalty $\sim 0.2\% \times f$ through different $\alpha$ scaling strategies
    \end{itemize}
\end{enumerate}

\textbf{Practical advantage of current design:}
\begin{itemize}
    \item \textbf{Automatic adaptation:} No need to manually tune $\alpha$ for different problem sizes
    \item \textbf{Learning flexibility:} RL can discover non-linear relationships between conflict and optimal penalty
    \item \textbf{Robustness:} Works well across diverse benchmark functions (F1-F6 validated)
\end{itemize}

\textbf{Relationship with Step 1 penalty:}
\begin{center}
\begin{tabular}{|l|c|c|}
\hline
\textbf{Aspect} & \textbf{Step 1 ($h_i$)} & \textbf{Step 3 (conflict)} \\
\hline
Purpose & Optimization guidance & Monitoring/Gating metric \\
Function & $h_i = f_i + \alpha \cdot \text{conflict}_i$ & $\text{conflict}_i$ (same formula) \\
Formula & $\sum_d \frac{|x_{i,d} - x_{i,\text{con}}^d|}{\text{range}}$ & $\sum_d \frac{|x_{i,\text{best}}^d - x_{i,\text{con}}^d|}{\text{range}}$ \\
Norm type & L1 (absolute) & L1 (absolute) \\
Normalization & Per-dimension (by range) & Per-dimension (by range) \\
Averaging & No (sum over dimensions) & No (sum over dimensions) \\
Weight & $\alpha(t)$ (adaptive) & 1 (unweighted) \\
Applied to & All $M$ particles & Best particle only \\
Range & $[0, |O_i|]$ & $[0, |O_i|]$ \\
\hline
\end{tabular}
\end{center}

\textbf{Key insight:} Step 3's conflict measure is \textbf{exactly the same} as the penalty term in Step 1, just evaluated on the best particle instead of all particles. This ensures consistency between optimization guidance and communication decisions.

\subsection{Step 4: Communication Index Gating (CI $>$ Threshold?)}

\textbf{Step 4.1: Compute raw Conflict Index}

Using L1-norm sum (NOT averaged), SAME formula as Step 3:
\begin{equation}
\text{CI}_i^{\text{raw}} = \sum_{d \in O_i} \frac{|x_{i,\text{best}}^d - x_{i,\text{con}}^d|}{x_{\max} - x_{\min}}
\end{equation}

\textbf{Important:} $\text{CI}_i^{\text{raw}} \in [0, |O_i|]$ scales with overlap size.

\textbf{Step 4.2: Exponential smoothing to reduce transient fluctuations}

To mitigate the impact of momentary noise, apply exponential moving average:
\begin{equation}
\text{CI}_i^{\text{smooth}} = \beta \cdot \text{CI}_i^{\text{smooth}} + (1 - \beta) \cdot \text{CI}_i^{\text{raw}}
\end{equation}
where $\beta = 0.7$ is the smoothing coefficient (70\% history + 30\% current).

\textbf{Step 4.3: Three-tier communication gating mechanism}

Communication is triggered only when ALL three conditions are satisfied:

\textbf{Condition 1: Minimum interval constraint}
\begin{equation}
\text{iter}_{\text{current}} - \text{iter}_{\text{last\_comm}} \geq T_{\min}
\end{equation}
Prevents excessive communication-induced oscillations. Default: $T_{\min} = 10$ iterations.

\textbf{Condition 2: Adaptive threshold test}
\begin{equation}
\text{CI}_i^{\text{smooth}} > \tau(t)
\end{equation}
where $\tau(t)$ is dynamically adjusted:
\begin{itemize}
    \item When conflicts are frequent: $\tau \downarrow$ (lower threshold, increase communication)
    \item When conflicts are sparse or communication benefit is low: $\tau \uparrow$ (raise threshold, save overhead)
    \item Specific adjustment rule:
    \begin{equation}
    \tau(t) = \begin{cases}
    0.5 \cdot \tau_{\text{base}} & \text{if } \text{CI}_i^{\text{smooth}} > 0.7 \text{ (high conflict)} \\
    2.0 \cdot \tau_{\text{base}} & \text{if } \text{CI}_i^{\text{smooth}} < 0.3 \text{ (low conflict)} \\
    \tau_{\text{base}} & \text{otherwise}
    \end{cases}
    \end{equation}
    where $\tau_{\text{base}} = 0.02$.
\end{itemize}

\textbf{Condition 3: Phase-based frequency sampling}

Use probabilistic gating to avoid deterministic over-communication:
\begin{equation}
\text{uniform\_random}(0,1) < f_{\text{phase}}(\text{iter}, \Delta f)
\end{equation}
where $\text{uniform\_random}(0,1)$ generates a random number in $[0,1]$, and $f_{\text{phase}}$ is the phase-dependent communication probability.

\textbf{Phase determination:} The optimization phase is determined by both iteration count and convergence rate:
\begin{equation}
\text{Phase}(\text{iter}, \Delta f) = \begin{cases}
\text{LATE} & \text{if } (\text{iter} > 30) \lor (\Delta f < 0.001) \\
\text{MID} & \text{if } (\text{iter} > 15) \lor (\Delta f < 0.01) \\
\text{EARLY} & \text{otherwise}
\end{cases}
\end{equation}
where $\Delta f$ is the fitness improvement rate (relative change per iteration).

\textbf{Frequency mapping:} Each phase has a base communication probability:
\begin{equation}
f_{\text{phase}} = \begin{cases}
0.6 & \text{if Phase = EARLY (aggressive exploration)} \\
0.3 & \text{if Phase = MID (balanced exploration-exploitation)} \\
0.15 & \text{if Phase = LATE (fine-tuning convergence)}
\end{cases}
\end{equation}

\textbf{Randomness rationale:} The probabilistic mechanism prevents:
\begin{itemize}
    \item Over-reaction to borderline CI values near the threshold
    \item Synchronous oscillations among multiple agents
    \item Excessive communication overhead in marginal conflict scenarios
\end{itemize}

\textbf{Final decision logic:}
\begin{align}
\text{should\_comm}_i &= (\text{iter} - \text{last\_comm} \geq T_{\min}) \land (\text{CI}_i^{\text{smooth}} > \tau(t)) \land (\text{rand}() < f_{\text{phase}}) \\
\text{global\_comm} &= \text{MPI\_Allreduce}(\text{should\_comm}_i, \text{MAX})
\end{align}

\subsection{Step 5: Exchange Solutions via MPI}

Agent $i$ communicates with all its neighbors in $\mathcal{N}_i$ (the set of neighboring agents sharing overlapping variables):

\textbf{Communication pattern:} For each neighbor $j \in \mathcal{N}_i$, perform bidirectional point-to-point exchange:
\begin{equation}
\begin{aligned}
&\text{Agent } i: \begin{cases}
\text{MPI\_Isend}(x_i^{\text{best}}, \text{destination}=j) & \text{(send own solution to } j\text{)} \\
\text{MPI\_Recv}(x_j^{\text{best}}, \text{source}=j) & \text{(receive } j\text{'s solution)}
\end{cases} \\
&\text{Agent } j: \begin{cases}
\text{MPI\_Isend}(x_j^{\text{best}}, \text{destination}=i) & \text{(send own solution to } i\text{)} \\
\text{MPI\_Recv}(x_i^{\text{best}}, \text{source}=i) & \text{(receive } i\text{'s solution)}
\end{cases}
\end{aligned}
\end{equation}

\textbf{Implementation note:}
\begin{itemize}
    \item Each agent loops through its neighbor set: \texttt{for (int rank : commu\_object)}
    \item Communication is symmetric: if $i$ sends to $j$, then $j$ sends to $i$ simultaneously
    \item Uses non-blocking \texttt{MPI\_Isend} to avoid deadlock, followed by blocking \texttt{MPI\_Recv}
    \item Each agent typically has 2-4 neighbors depending on variable overlap structure
\end{itemize}

\subsection{Step 6: Nash Bargaining Competition}

\textbf{Notation convention (critical):}
\begin{itemize}
    \item $x_{i,\text{best}}^d$: the $d$-th dimension of agent $i$'s best solution (subscript = agent, superscript = dimension)
    \item $f_{a,b}^d$: \textbf{agent $a$'s solution evaluated by agent $b$'s objective function}
    \begin{itemize}
        \item First subscript: whose solution (source of decision variable values)
        \item Second subscript: whose objective function (evaluator's perspective)
    \end{itemize}
\end{itemize}

\textbf{Step 6.1: Evaluate cross-impact fitness}

For each overlapping dimension $d$ shared between agents $i$ and $j$, compute four fitness values representing mutual impact:

\begin{align}
f_{i,i}^d &= f_i(x_{i,\text{best}}) \quad \text{(agent } i\text{'s solution, evaluated by } i\text{'s objective)} \\
f_{i,j}^d &= f_j(x_{i,\text{best}}) \quad \text{(agent } i\text{'s solution, evaluated by } j\text{'s objective)} \\
f_{j,i}^d &= f_i(x_{j,\text{best}}) \quad \text{(agent } j\text{'s solution, evaluated by } i\text{'s objective)} \\
f_{j,j}^d &= f_j(x_{j,\text{best}}) \quad \text{(agent } j\text{'s solution, evaluated by } j\text{'s objective)}
\end{align}

\textbf{Interpretation:}
\begin{itemize}
    \item $f_{i,j}^d$: How friendly is agent $i$'s solution to agent $j$'s sub-problem? (cooperation measure)
    \item $f_{j,i}^d$: How friendly is agent $j$'s solution to agent $i$'s sub-problem?
\end{itemize}

\textbf{Step 6.2: Nash bargaining decision}

Apply Nash bargaining solution to determine consensus value for dimension $d$:
\begin{equation}
x_{i,\text{consensus}}^d = \begin{cases}
x_{j,\text{best}}^d & \text{if } f_{i,i}^d + f_{j,i}^d > f_{i,j}^d + f_{j,j}^d \quad \text{(adopt } j\text{'s value)} \\
x_{i,\text{best}}^d & \text{otherwise} \quad \text{(keep } i\text{'s value)}
\end{cases}
\end{equation}

\textbf{Economic interpretation:}
\begin{itemize}
    \item \textbf{Left side} $f_{i,i}^d + f_{j,i}^d$: Joint utility when adopting agent $i$'s value
    \begin{itemize}
        \item $f_{i,i}^d$: benefit to agent $i$ from its own solution
        \item $f_{j,i}^d$: benefit to agent $i$ from agent $j$'s perspective (cross-impact)
    \end{itemize}
    \item \textbf{Right side} $f_{i,j}^d + f_{j,j}^d$: Joint utility when adopting agent $j$'s value
    \begin{itemize}
        \item $f_{i,j}^d$: benefit to agent $j$ from agent $i$'s perspective (cross-impact)
        \item $f_{j,j}^d$: benefit to agent $j$ from its own solution
    \end{itemize}
    \item \textbf{Decision rule}: Choose the value maximizing joint welfare (cooperative Pareto improvement)
\end{itemize}

\subsection{Step 7: Gradient-based Conflict Detection}
Compute numerical gradients using finite difference:
\begin{align}
g_i^{+}(d) &= f_i(x_1, \ldots, x_d + \delta, \ldots, x_D) - f_i(x) \\
g_i^{-}(d) &= f_i(x_1, \ldots, x_d - \delta, \ldots, x_D) - f_i(x)
\end{align}
Reject consensus if gradients agree (no conflict):
\begin{equation}
\text{Reject if: } [\text{sign}(g_i^+) = \text{sign}(g_j^+)] \land [\text{sign}(g_i^-) = \text{sign}(g_j^-)]
\end{equation}

\subsection{Step 8: Skip Negotiation}
If communication is not triggered (CI $\leq$ threshold and not forced interval):
\begin{equation}
x_{i,\text{con}} = x_{i,\text{best}} \quad \text{(directly use local best from Step 2, skip Steps 5-7)}
\end{equation}

\textbf{Logic:} When conflict is low, negotiation is unnecessary. Agent $i$ directly adopts its own local best solution without communication overhead.

\subsection{Step 9: Set Global Best}
Update global best solution for each agent based on whether negotiation occurred:
\begin{equation}
x_{i,\text{con}} = \begin{cases}
x_{i,\text{consensus}} & \text{if negotiation occurred (Step 6)} \\
x_{i,\text{best}} & \text{if negotiation skipped (Step 8)}
\end{cases}
\end{equation}

\textbf{Purpose:} Set the consensus reference $x_{i,\text{con}}$ (variable name: \texttt{globalBest}) that will be used in:
\begin{itemize}
    \item Penalty term calculation: $\sum_{d \in O_i}(x_{i,d} - x_{i,\text{con}}^d)^2$ in next iteration's Step 1
    \item Conflict calculation: compare $x_{i,\text{best}}^d$ with $x_{i,\text{con}}^d$ in Step 3
\end{itemize}

\textbf{Iteration flow:}
\begin{itemize}
    \item Iteration $t-1$: produces $x_{i,\text{con}}^{(t-1)}$ (stored in memory)
    \item Iteration $t$: Step 1 uses $x_{i,\text{con}}^{(t-1)}$ as penalty reference $\to$ Step 9 updates to $x_{i,\text{con}}^{(t)}$
\end{itemize}

\subsection{Step 10: Re-evaluate Population}
Re-evaluate all particles with updated consensus reference:
\begin{equation}
\forall p \in \text{Population}_i: \quad h_i(p) = f_i(p) + \alpha_i(t) \cdot \text{conflict}_i(p)
\end{equation}

where
\begin{equation}
\text{conflict}_i(p) = \sum_{d \in O_i} \frac{|p_{i,d} - x_{i,\text{con}}^d|}{x_{\max} - x_{\min}}
\end{equation}

\textbf{Purpose:} After negotiation updates $x_{i,\text{con}}$ (Step 9), recalculate penalized fitness $h_i$ for all particles using the new consensus reference.

\textbf{Notation clarification:}
\begin{itemize}
    \item $p \in \text{Population}_i$: A particle (solution vector) in agent $i$'s population
    \item $p_{i,d}$: The value of dimension $d$ in particle $p$ (consistent with Step 1's notation $x_{i,d}$)
    \item $\alpha_i(t)$: Time-varying penalty coefficient (same as Step 1)
    \item $\text{conflict}_i(p)$: Conflict measure for particle $p$ (same formula as Step 1 and Step 3)
    \item This formula has the \textbf{exact same structure} as Step 1, just applied to all particles after consensus update
    \item \textbf{No averaging:} Sum over all overlapping dimensions, consistent with Steps 1 and 3
\end{itemize}

\subsection{Step 11: MPI Allreduce - Global Fitness}
Aggregate fitness values across all agents to compute global metrics:
\begin{align}
f^{\text{pure}} &= \frac{1}{N}\sum_{i=1}^N f_i(x_{i,\text{con}}) = \text{MPI\_Allreduce}(\text{SUM}) / N \\
h^{\text{penalty}} &= \frac{1}{N}\sum_{i=1}^N h_i(x_{i,\text{con}}) = \text{MPI\_Allreduce}(\text{SUM}) / N \\
\text{penalty} &= h^{\text{penalty}} - f^{\text{pure}}
\end{align}

\textbf{Understanding the penalty term:}

Since $h_i(x_{i,\text{con}}) = f_i(x_{i,\text{con}}) + \alpha_i(t) \cdot \text{conflict}_i(x_{i,\text{con}})$, we have:
\begin{equation}
\text{penalty} = \frac{1}{N}\sum_{i=1}^N \left[\alpha_i(t) \cdot \text{conflict}_i(x_{i,\text{con}})\right]
\end{equation}

where
\begin{equation}
\text{conflict}_i(x_{i,\text{con}}) = \sum_{d \in O_i} \frac{|x_{i,\text{con}}^d - x_{i,\text{con}}^d|}{x_{\max} - x_{\min}} = 0 \quad \text{(at perfect consensus)}
\end{equation}

\textbf{Important clarification on conflict range:}
\begin{itemize}
    \item In the current implementation: $\text{conflict}_i \in [0, |O_i|]$ (sum, not average)
    \item For typical problems: $|O_i| = 100 \sim 300$ overlapping dimensions
    \item At initialization: $\text{conflict}_i \approx 0.5 \times |O_i|$ (random deviations)
    \item Example: $|O_i| = 200$, average deviation $= 0.5 \Rightarrow \text{conflict}_i \approx 100$
    \item The RL-adjusted $\alpha$ compensates for this scaling to maintain effective penalty strength
\end{itemize}

\textbf{Key distinction:}
\begin{itemize}
    \item $f^{\text{pure}} = \frac{1}{N}\sum_{i=1}^N f_i(x_{i,\text{con}})$: Average of \textbf{pure} local fitness (no penalty)
    \item $h^{\text{penalty}} = \frac{1}{N}\sum_{i=1}^N h_i(x_{i,\text{con}})$: Average of \textbf{penalized} local fitness
    \item $\text{penalty} = \frac{1}{N}\sum_{i=1}^N \alpha_i(t) \cdot \text{conflict}_i(x_{i,\text{con}})$: Average weighted conflict
    \item $F(x) = \sum_{i=1}^N f_i(x_i)$: Global objective (without averaging)
    \item At convergence: $\text{conflict}_i \to 0$, thus $\text{penalty} \to 0$, and $h^{\text{penalty}} \to f^{\text{pure}}$
\end{itemize}

\textbf{Typical penalty magnitude analysis:}
\begin{itemize}
    \item Early iterations: $\text{conflict}_i \approx 50 \sim 150$ (depending on $|O_i|$ and initialization)
    \item With $\alpha \approx 0.001$ (RL-adjusted): $\text{penalty} \approx 0.05 \sim 0.15$
    \item For $f^{\text{pure}} \approx 10^5$: $\text{penalty}/f \approx 0.0005\% \sim 0.0015\%$ (very small, as designed)
    \item The RL mechanism dynamically tunes $\alpha$ to balance exploration and convergence
\end{itemize>

\textbf{Important note:} When evaluating $h_i(x_{i,\text{con}})$, the conflict term compares $x_{i,\text{con}}$ with itself (since $x_{i,\text{con}}$ is the consensus reference). At perfect consensus across all agents, this conflict becomes zero. However, during optimization, different agents may have slightly different $x_{i,\text{con}}$ values, leading to non-zero penalties.

\subsection{Step 12: Calculate Improvement}
Measure relative fitness improvement:
\begin{equation}
\text{improvement}_t = \frac{f_{t-1}^{\text{pure}} - f_t^{\text{pure}}}{|f_{t-1}^{\text{pure}}|}
\end{equation}

\subsection{Step 13: Compute Reward}
RL reward signal uses Scheme 5 (dimensionless penalty-efficiency reward):
\begin{equation}
r_t = \tanh\!\left(-\Delta\text{conflict}_{\text{rate}} \cdot \Delta f_{\text{pure,rate}} \cdot 100\right)
\end{equation}
where
\begin{align}
\Delta\text{conflict}_{\text{rate}} &= \frac{\text{conflict}_t - \text{conflict}_{t-1}}{\text{conflict}_{t-1} + 10^{-10}}, \\
\Delta f_{\text{pure,rate}} &= \frac{f^{\text{pure}}_{t-1} - f^{\text{pure}}_t}{|f^{\text{pure}}_{t-1}| + 10^{-10}} .
\end{align}
\textbf{Interpretation:} reward is positive when conflict decreases and pure fitness improves simultaneously.

\subsection{Step 14: RL Update Penalty Coefficient $\alpha$}
\textbf{Key insight:} This step uses the reward $r_t$ from Step 13 to train the RL policy network, enabling adaptive penalty adjustment.

\subsubsection*{14.1 State Construction}
Construct the environment state from current conflict and historical trend:
\begin{equation}
s_t = [S_1, S_2] \in \mathbb{R}^2
\end{equation}

where:
\begin{align}
S_1 &= \text{conflict}_i(x_{i,\text{con}}) = \sum_{d \in O_i} \frac{|x_{i,\text{best}}^d - x_{i,\text{con}}^d|}{x_{\max} - x_{\min}} \quad \text{(current conflict, from Step 3)} \\
S_2 &= \text{conflict}_i^{(t)} - \text{conflict}_i^{(t-1)} \quad \text{(first-order difference, conflict rate of change)}
\end{align}

\textbf{Implementation details:}
\begin{itemize}
    \item $S_1$ is the \textbf{raw conflict value} computed in Step 3 (not the smoothed CI from Step 4)
    \item $S_2$ is a simple \textbf{one-step difference} (not a complex trend metric)
    \item Conflict history is maintained in a sliding window of size 10
    \item If no history exists (first iteration), $S_2 = 0$
\end{itemize}

\textbf{Numerical scaling:}
\begin{itemize}
    \item $S_1 \in [0, |O_i|]$ (sum of normalized deviations across overlapping dimensions)
    \item For typical problems: $S_1 \in [0, 300]$ depending on overlap size and deviation
    \item $S_2 \in [-|O_i|, |O_i|]$ (maximum possible change in conflict)
    \item Example: $S_1 = 2.0$ (low conflict), $S_2 = -0.5$ (conflict decreasing)
\end{itemize}

\textbf{Connection to Step 3 and Step 4:}
\begin{itemize}
    \item Step 3 computes \textbf{raw conflict} $\text{conflict}_i$ (instantaneous value)
    \item Step 4 uses \textbf{smoothed CI} $\text{CI}_i^{\text{smooth}}$ for gating decision (stable threshold comparison)
    \item Step 14.1 uses \textbf{raw conflict as $S_1$} for RL state (fast feedback for learning)
\end{itemize}

\textbf{Critical distinction: $S_1$ vs $\text{CI}_i^{\text{smooth}}$}

\begin{center}
\begin{tabular}{|l|c|c|}
\hline
\textbf{Aspect} & \textbf{$S_1$ (RL State)} & \textbf{$\text{CI}_i^{\text{smooth}}$ (Gating)} \\
\hline
\textbf{Formula} & $S_1 = \text{conflict}_i^{\text{raw}}$ & $\text{CI}^{\text{smooth}} = 0.7 \cdot \text{CI}^{\text{smooth}} + 0.3 \cdot \text{conflict}_i^{\text{raw}}$ \\
\hline
\textbf{Smoothing} & ❌ No (raw value) & ✅ Yes (EMA with $\beta=0.7$) \\
\hline
\textbf{Response speed} & ⚡ Instant & 🐌 Delayed (3-5 iterations lag) \\
\hline
\textbf{Noise sensitivity} & High (reflects transient spikes) & Low (filters out noise) \\
\hline
\textbf{Usage} & RL training feedback & Communication threshold test \\
\hline
\textbf{Purpose} & Learn adaptive policy & Stable gate control \\
\hline
\textbf{Typical value} & $0.001 \sim 0.05$ & $0.001 \sim 0.05$ (same range, smoother) \\
\hline
\end{tabular}
\end{center}

\textbf{Why use raw conflict ($S_1$) for RL instead of smoothed CI?}

\begin{itemize}
    \item \textbf{RL needs immediate feedback}: The policy network must observe the true consequence of its actions ($\alpha$ adjustment) without delay. Smoothing would corrupt the causal relationship between action and reward.
    \item \textbf{RL can handle noise}: Neural networks naturally smooth noisy signals through weight averaging across multiple experiences. External smoothing is redundant and harmful.
    \item \textbf{Example}: If RL increases $\alpha$ at $t=10$, it should see the conflict change at $t=11$, not a smoothed average of $t=7$ to $t=11$.
\end{itemize}

\textbf{Numerical example comparing $S_1$ and $\text{CI}^{\text{smooth}}$:}

\begin{center}
\begin{tabular}{|c|c|c|c|}
\hline
\textbf{Iteration} & \textbf{Raw conflict} & \textbf{$S_1$ (RL State)} & \textbf{$\text{CI}^{\text{smooth}}$ (Gate)} \\
\hline
1 & 0.0500 & 0.0500 & 0.0500 (init) \\
2 & 0.0450 & 0.0450 & $0.7 \times 0.0500 + 0.3 \times 0.0450 = 0.0485$ \\
3 & 0.0600 (spike) & \textbf{0.0600} ⚡ & $0.7 \times 0.0485 + 0.3 \times 0.0600 = 0.0520$ 🐌 \\
4 & 0.0400 & 0.0400 & $0.7 \times 0.0520 + 0.3 \times 0.0400 = 0.0484$ \\
\hline
\end{tabular}
\end{center}

\textbf{Observation from example:}
\begin{itemize}
    \item At $t=3$: $S_1$ immediately jumps to 0.0600 (RL sees the spike instantly)
    \item At $t=3$: $\text{CI}^{\text{smooth}}$ only rises to 0.0520 (gate sees a dampened signal)
    \item At $t=4$: $S_1$ drops to 0.0400 (RL sees recovery), but $\text{CI}^{\text{smooth}}$ is still high at 0.0484
    \item This lag in $\text{CI}^{\text{smooth}}$ is desirable for gating (avoids oscillation) but harmful for RL learning
\end{itemize}

\subsubsection*{14.2 Action Selection (Policy Network)}
The policy network $\pi_\theta$ maps state to weight adjustment action:
\begin{align}
a_t &= \pi_\theta(s_t) = \tanh(W \cdot s_t + b) \\
\text{ratio}_{t+1} &= \text{clip}(\text{ratio}_t + 0.1 \cdot a_t, 0.1, 2.0), \\
\alpha_{t+1} &= \text{base\_alpha}_t \cdot \text{ratio}_{t+1}, \quad \text{base\_alpha}_t = |f_t^{\text{pure}}|/512
\end{align}

\textbf{Network architecture and dimensions:}

\begin{itemize}
    \item \textbf{Input}: $s_t = [S_1, S_2]^T \in \mathbb{R}^{2 \times 1}$ (column vector, 2-dimensional state)
    \item \textbf{Weight matrix}: $W \in \mathbb{R}^{1 \times 2}$ (1 row, 2 columns)
    \item \textbf{Bias}: $b \in \mathbb{R}$ (scalar)
    \item \textbf{Output}: $a_t \in \mathbb{R}$ (scalar action)
\end{itemize}

\textbf{Matrix multiplication detailed:}
\begin{equation}
W \cdot s_t = \begin{bmatrix} w_1 & w_2 \end{bmatrix} \cdot \begin{bmatrix} S_1 \\ S_2 \end{bmatrix} = w_1 \cdot S_1 + w_2 \cdot S_2 \in \mathbb{R}
\end{equation}

where $w_1, w_2$ are the learnable weights corresponding to conflict and trend features.

\textbf{Complete forward pass:}
\begin{align}
z &= W \cdot s_t + b = w_1 S_1 + w_2 S_2 + b \quad \text{(linear combination)} \\
a_t &= \tanh(z) = \frac{e^z - e^{-z}}{e^z + e^{-z}} \in [-1, 1] \quad \text{(nonlinear activation)}
\end{align}

\textbf{Dimension check:}
\begin{align}
W \cdot s_t: & \quad \underbrace{\mathbb{R}^{1 \times 2}}_{W} \times \underbrace{\mathbb{R}^{2 \times 1}}_{s_t} = \mathbb{R}^{1 \times 1} = \mathbb{R} \quad \checkmark \\
z + b: & \quad \underbrace{\mathbb{R}}_{W \cdot s_t} + \underbrace{\mathbb{R}}_{b} = \mathbb{R} \quad \checkmark \\
\tanh(z): & \quad \mathbb{R} \rightarrow \mathbb{R} \quad \checkmark
\end{align}

\textbf{Parameter explanation:}
\begin{itemize}
    \item $W \in \mathbb{R}^{1 \times 2}$ means: \textbf{1 output neuron, 2 input features}
    \item Expanded: $W = [w_1, w_2]$ where $w_1$ weights $S_1$ (conflict), $w_2$ weights $S_2$ (trend)
    \item Total trainable parameters: $\theta = \{w_1, w_2, b\}$ (3 parameters)
    \item $\tanh(\cdot)$ bounds output to $[-1, 1]$ (ensures stable weight updates)
    \item $0.01$ is the step size for weight adjustment: $\Delta\alpha = a_t \times 0.01$
    \item $[0.01, 1.0]$ is the valid range for $\alpha$ (hard constraints)
\end{itemize}

\textbf{Numerical example:}
\begin{align*}
\text{Assume: } & W = [2.5, -1.8], \quad b = 0.3 \\
\text{State: } & s_t = [0.01, -0.003]^T \quad \text{(low conflict, decreasing)} \\
\text{Compute: } & z = 2.5 \times 0.01 + (-1.8) \times (-0.003) + 0.3 = 0.0304 \\
& a_t = \tanh(0.0304) \approx 0.0304 \\
& \Delta\alpha = 0.0304 \times 0.01 = 0.000304 \\
\text{Result: } & \alpha_{t+1} = \alpha_t + 0.000304 \quad \text{(slight increase in penalty)}
\end{align*}

\subsubsection*{14.3 Experience Storage}
\textbf{Critical step:} Store the decision experience using reward from Step 13:
\begin{equation}
\text{Experience: } \mathcal{E}_t = (s_t, a_t, r_t) \rightarrow \text{ReplayBuffer}
\end{equation}

\textbf{Priority assignment:}
\begin{equation}
p_t = |r_t| \quad \text{(priority value stored with experience)}
\end{equation}

where $p_t$ is the \textbf{priority value} (not probability) assigned to experience $\mathcal{E}_t$.

\textbf{Key points:}
\begin{itemize}
    \item Priority value $p_t = |r_t|$ is stored directly (e.g., if $r_t = -18.4$, then $p_t = 18.4$)
    \item Higher $|r_t|$ (large reward magnitude) → higher priority value
    \item This priority value will be used in Step 14.4 to determine sampling probability
    \item Both positive and negative rewards contribute equally (absolute value)
\end{itemize}

\subsubsection*{14.4 Policy Network Update (Supervised Learning)}
When replay buffer has $\geq 32$ experiences, perform batch update:
\begin{align}
a_{\text{target}} &= a_t + 0.1 \times r_t \quad \text{(reward-corrected action)} \\
\mathcal{L}(\theta) &= (a_{\text{target}} - \pi_\theta(s_t))^2 \quad \text{(MSE loss)} \\
\theta &\leftarrow \theta - \eta \nabla_\theta \mathcal{L}(\theta) \quad \text{(gradient descent)}
\end{align}

\textbf{Learning rate $\eta$:}
\begin{equation}
\eta = 0.01 \quad \text{(fixed learning rate)}
\end{equation}

\textbf{Parameter update details:}

\textbf{What is $\theta$? Understanding the parameter set}

The symbol $\theta$ represents the \textbf{set of all trainable parameters} in the policy network $\pi_\theta$. 

\textbf{Why $\theta = \{w_1, w_2, b\}$?}

The policy network has a simple structure with \textbf{exactly 3 trainable parameters}:

\begin{enumerate}
    \item \textbf{Weight $w_1$}: Connects input $S_1$ (conflict) to output $a_t$
    \item \textbf{Weight $w_2$}: Connects input $S_2$ (trend) to output $a_t$
    \item \textbf{Bias $b$}: Adds a constant offset to the weighted sum
\end{enumerate}

\textbf{Network structure:}
\begin{equation}
\pi_\theta(s_t) = \tanh(\underbrace{w_1 \cdot S_1 + w_2 \cdot S_2}_{\text{weighted sum}} + \underbrace{b}_{\text{bias}})
\end{equation}

\textbf{Parameter count derivation:}
\begin{itemize}
    \item \textbf{Weight matrix $W \in \mathbb{R}^{1 \times 2}$}: Contains 2 elements $\{w_1, w_2\}$
    \item \textbf{Bias $b \in \mathbb{R}$}: Contains 1 element $\{b\}$
    \item \textbf{Total}: $|\theta| = 2 + 1 = 3$ parameters
\end{itemize}

\textbf{Notation explanation:}
\begin{itemize}
    \item $\theta$ = \textbf{parameter set} (all trainable variables)
    \item $\theta_i$ = \textbf{individual parameter} (one element from the set)
    \item $\theta_i \in \{w_1, w_2, b\}$ means: "for each parameter $i$ that is either $w_1$, $w_2$, or $b$"
\end{itemize}

\textbf{Why only these 3 parameters?}
\begin{itemize}
    \item The network has \textbf{no hidden layers} (direct input-to-output mapping)
    \item Input dimension = 2 ($S_1, S_2$), Output dimension = 1 ($a_t$)
    \item Required weights: $2 \times 1 = 2$ (one weight per input-output connection)
    \item Required bias: $1$ (one bias per output neuron)
    \item Total: $2 + 1 = 3$ parameters
\end{itemize}

\textbf{Gradient computation for each parameter:}

For each parameter $\theta_i \in \{w_1, w_2, b\}$:
\begin{align}
\nabla_{\theta_i} \mathcal{L} &= \frac{\partial}{\partial \theta_i} (a_{\text{target}} - \pi_\theta(s_t))^2 \\
&= 2(a_{\text{target}} - \pi_\theta(s_t)) \cdot \frac{\partial \pi_\theta(s_t)}{\partial \theta_i}
\end{align}

This means we compute the gradient \textbf{separately} for each of the 3 parameters:
\begin{align}
\nabla_{w_1} \mathcal{L} &= 2(a_{\text{target}} - a_t) \cdot \frac{\partial \pi_\theta(s_t)}{\partial w_1} \\
\nabla_{w_2} \mathcal{L} &= 2(a_{\text{target}} - a_t) \cdot \frac{\partial \pi_\theta(s_t)}{\partial w_2} \\
\nabla_{b} \mathcal{L} &= 2(a_{\text{target}} - a_t) \cdot \frac{\partial \pi_\theta(s_t)}{\partial b}
\end{align}

\textbf{Complete update equations:}
\begin{align}
\frac{\partial \pi_\theta(s_t)}{\partial w_1} &= (1 - \tanh^2(z)) \cdot S_1 \quad \text{where } z = w_1 S_1 + w_2 S_2 + b \\
\frac{\partial \pi_\theta(s_t)}{\partial w_2} &= (1 - \tanh^2(z)) \cdot S_2 \\
\frac{\partial \pi_\theta(s_t)}{\partial b} &= (1 - \tanh^2(z))
\end{align}

Therefore:
\begin{align}
w_1 &\leftarrow w_1 - \eta \cdot 2(a_{\text{target}} - a_t) \cdot (1 - a_t^2) \cdot S_1 \\
w_2 &\leftarrow w_2 - \eta \cdot 2(a_{\text{target}} - a_t) \cdot (1 - a_t^2) \cdot S_2 \\
b &\leftarrow b - \eta \cdot 2(a_{\text{target}} - a_t) \cdot (1 - a_t^2)
\end{align}

where $a_t = \pi_\theta(s_t) = \tanh(z)$ is the current network output.

\textbf{Why $\eta = 0.01$?}
\begin{itemize}
    \item \textbf{Small enough}: Prevents overshooting optimal weights (stable convergence)
    \item \textbf{Large enough}: Allows meaningful updates within reasonable training time
    \item \textbf{Empirical choice}: Standard value for simple neural networks (typical range: $0.001 \sim 0.1$)
    \item \textbf{Fixed}: No adaptive scheduling (simplicity over sophistication)
\end{itemize}

\textbf{What if we use adaptive learning rate? Analysis of alternatives}

\textbf{Common adaptive methods:}

\begin{enumerate}
    \item \textbf{AdaGrad}: Accumulates squared gradients
    \begin{equation}
    \eta_t^{(i)} = \frac{\eta_0}{\sqrt{G_t^{(i)} + \epsilon}}, \quad G_t^{(i)} = \sum_{k=1}^{t} (\nabla_{\theta_i} \mathcal{L}_k)^2
    \end{equation}
    where $\eta_0$ is initial learning rate, $\epsilon = 10^{-8}$ prevents division by zero.
    
    \item \textbf{RMSprop}: Exponential moving average of squared gradients
    \begin{equation}
    \eta_t^{(i)} = \frac{\eta_0}{\sqrt{v_t^{(i)} + \epsilon}}, \quad v_t^{(i)} = \beta v_{t-1}^{(i)} + (1-\beta)(\nabla_{\theta_i} \mathcal{L}_t)^2
    \end{equation}
    where $\beta = 0.9$ is decay rate.
    
    \item \textbf{Adam}: Combines momentum and adaptive learning rate
    \begin{align}
    m_t^{(i)} &= \beta_1 m_{t-1}^{(i)} + (1-\beta_1)\nabla_{\theta_i} \mathcal{L}_t \quad \text{(momentum)} \\
    v_t^{(i)} &= \beta_2 v_{t-1}^{(i)} + (1-\beta_2)(\nabla_{\theta_i} \mathcal{L}_t)^2 \quad \text{(variance)} \\
    \hat{m}_t^{(i)} &= \frac{m_t^{(i)}}{1-\beta_1^t}, \quad \hat{v}_t^{(i)} = \frac{v_t^{(i)}}{1-\beta_2^t} \quad \text{(bias correction)} \\
    \eta_t^{(i)} &= \frac{\eta_0 \cdot \hat{m}_t^{(i)}}{\sqrt{\hat{v}_t^{(i)}} + \epsilon}
    \end{align}
    where $\beta_1 = 0.9$, $\beta_2 = 0.999$ are momentum decay rates.
\end{enumerate}

\textbf{Potential benefits for MACPO:}
\begin{itemize}
    \item \textbf{Faster convergence}: Large gradients → smaller learning rate (prevents overshooting), small gradients → larger learning rate (faster progress)
    \item \textbf{Per-parameter adaptation}: Different learning rates for $w_1$, $w_2$, $b$ based on their gradient magnitudes
    \item \textbf{Automatic tuning}: Reduces need for manual $\eta$ selection
\end{itemize}

\textbf{Potential drawbacks for MACPO:}
\begin{itemize}
    \item \textbf{Overfitting risk}: Adaptive methods can converge too quickly to local minima, especially with only 3 parameters
    \item \textbf{Hyperparameter sensitivity}: AdaGrad/RMSprop/Adam introduce new hyperparameters ($\beta_1, \beta_2, \epsilon$) that need tuning
    \item \textbf{Memory overhead}: Need to store momentum/variance for each parameter (3× more memory)
    \item \textbf{Computational overhead}: Additional operations per update (square root, division, exponential moving average)
    \item \textbf{Instability}: Adaptive methods can be unstable with sparse gradients (common in RL)
\end{itemize}

\textbf{Complexity analysis:}

\begin{center}
\begin{tabular}{|l|c|c|c|c|}
\hline
\textbf{Aspect} & \textbf{Fixed $\eta$} & \textbf{AdaGrad} & \textbf{RMSprop} & \textbf{Adam} \\
\hline
\textbf{Code lines} & 1 & +15 & +20 & +35 \\
\hline
\textbf{Memory per param} & 0 & +1 (G) & +1 (v) & +2 (m, v) \\
\hline
\textbf{Operations per update} & 3 mult & +3 (sqrt, div) & +4 (EMA, sqrt) & +8 (momentum, EMA, bias corr) \\
\hline
\textbf{Hyperparameters} & 1 ($\eta$) & 2 ($\eta_0, \epsilon$) & 3 ($\eta_0, \beta, \epsilon$) & 4 ($\eta_0, \beta_1, \beta_2, \epsilon$) \\
\hline
\textbf{Convergence speed} & Baseline & Faster (early) & Moderate & Fastest \\
\hline
\textbf{Stability} & High & Medium & Medium & Low (can diverge) \\
\hline
\end{tabular}
\end{center}

\textbf{For MACPO's 3-parameter network:}
\begin{itemize}
    \item \textbf{Memory increase}: AdaGrad: +3 floats, RMSprop: +3 floats, Adam: +6 floats (negligible for modern systems)
    \item \textbf{Computation increase}: ~2-3× more operations per update (still very fast, < 1$\mu$s)
    \item \textbf{Code complexity}: Moderate increase (15-35 lines), but adds maintenance burden
\end{itemize}

\textbf{Why MACPO chooses fixed $\eta$:}
\begin{itemize}
    \item \textbf{Simplicity}: Only 3 parameters, fixed $\eta = 0.01$ works well empirically
    \item \textbf{Stability}: RL training benefits from predictable, stable updates
    \item \textbf{Interpretability}: Fixed learning rate makes behavior easier to understand and debug
    \item \textbf{Minimal benefit}: Adaptive methods shine with large networks (1000+ parameters), not 3-parameter networks
    \item \textbf{Hyperparameter avoidance}: One less thing to tune = faster development
\end{itemize}

\textbf{When adaptive learning rate might help:}
\begin{itemize}
    \item If MACPO network is expanded to include hidden layers (10+ parameters)
    \item If training shows clear signs of slow convergence or oscillation
    \item If different parameters ($w_1$ vs $w_2$ vs $b$) have vastly different gradient scales
    \item If computational overhead is not a concern
\end{itemize}

\textbf{Recommendation:}
\begin{itemize}
    \item \textbf{Current design}: Keep fixed $\eta = 0.01$ ✅ (simple, stable, sufficient)
    \item \textbf{Future consideration}: If network grows to 10+ parameters, consider RMSprop or Adam
    \item \textbf{Experimental option}: Allow adaptive learning rate as an optional feature (configurable)
\end{itemize}

\textbf{Numerical example:}

Assume:
\begin{align*}
s_t &= [0.01, -0.003]^T \\
a_t &= \pi_\theta(s_t) = 0.03 \quad \text{(current output)} \\
a_{\text{target}} &= 0.05 \quad \text{(target after reward correction)} \\
\eta &= 0.01
\end{align*}

Update computation:
\begin{align*}
\text{Error} &= a_{\text{target}} - a_t = 0.05 - 0.03 = 0.02 \\
\text{tanh derivative} &= 1 - a_t^2 = 1 - 0.03^2 = 0.9991 \\
\text{Gradient factor} &= 2 \times 0.02 \times 0.9991 = 0.03996 \\
\Delta w_1 &= -\eta \times 0.03996 \times S_1 = -0.01 \times 0.03996 \times 0.01 = -0.000004 \\
\Delta w_2 &= -\eta \times 0.03996 \times S_2 = -0.01 \times 0.03996 \times (-0.003) = +0.0000012 \\
\Delta b &= -\eta \times 0.03996 = -0.01 \times 0.03996 = -0.0004
\end{align*}

\textbf{Observation:}
\begin{itemize}
    \item Weight updates are very small ($\sim 10^{-6}$ for $w_1, w_2$)
    \item Bias update is larger ($\sim 10^{-4}$) but still conservative
    \item This ensures stable learning without oscillations
\end{itemize}

\textbf{Priority-based sampling:}

When sampling 32 experiences from the replay buffer, the sampling probability is \textbf{proportional to} the priority value:
\begin{equation}
P(\text{select experience } i) = \frac{p_i}{\sum_{j=1}^{N} p_j} = \frac{|r_i|}{\sum_{j=1}^{N} |r_j|} \quad \text{(sampling probability)}
\end{equation}

where $N$ is the current buffer size ($N \geq 32$).

\textbf{Key distinction:}
\begin{itemize}
    \item \textbf{Step 14.3}: Priority \textbf{value} $p_t = |r_t|$ (stored directly, e.g., $p_t = 18.4$)
    \item \textbf{Step 14.4}: Sampling \textbf{probability} $P(i) \propto |r_t|$ (normalized, e.g., $P(i) = 0.15$)
    \item The relationship: $P(i) = p_i / \sum_j p_j$ (probability is normalized priority)
\end{itemize}

\textbf{Numerical example:}

Assume replay buffer contains 3 experiences:
\begin{align*}
\mathcal{E}_1: & \quad r_1 = -18.4 \Rightarrow p_1 = 18.4 \\
\mathcal{E}_2: & \quad r_2 = -13.8 \Rightarrow p_2 = 13.8 \\
\mathcal{E}_3: & \quad r_3 = -9.2 \Rightarrow p_3 = 9.2 \\
\text{Total: } & \quad \sum p_j = 18.4 + 13.8 + 9.2 = 41.4
\end{align*}

Sampling probabilities:
\begin{align*}
P(1) &= \frac{18.4}{41.4} = 0.444 \quad \text{(44.4\%, highest priority)} \\
P(2) &= \frac{13.8}{41.4} = 0.333 \quad \text{(33.3\%)} \\
P(3) &= \frac{9.2}{41.4} = 0.222 \quad \text{(22.2\%, lowest priority)}
\end{align*}

\textbf{Training mechanism explained:}
\begin{itemize}
    \item \textbf{Positive reward} ($r_t > 0$): The target action $a_{\text{target}}$ increases, encouraging the network to output larger weight adjustments in similar states (more aggressive penalty)
    \item \textbf{Negative reward} ($r_t < 0$): The target action decreases, teaching the network to reduce weight adjustments (gentler penalty)
    \item \textbf{Batch sampling}: 32 experiences are sampled with probability $P(i) \propto |r_t|$, so high-impact decisions (large $|r_t|$) are sampled more frequently and learned more intensively
\end{itemize}

\textbf{Example training sequence:}
\begin{align*}
t=1: \quad & s_1 = [0.01, -0.5], \; a_1 = 0.02, \; r_1 = -18.4 \\
           & a_{\text{target}} = 0.02 + 0.1 \times (-18.4) = -1.82 \\
           & \text{Network learns: In high-conflict state, reduce weight aggressively} \\[5pt]
t=10: \quad & s_{10} = [0.002, -0.8], \; a_{10} = -0.01, \; r_{10} = -13.8 \\
            & a_{\text{target}} = -0.01 + 0.1 \times (-13.8) = -1.39 \\
            & \text{Network learns: Continue reducing weight when conflict decreases}
\end{align*}

\textbf{Connection to Step 13:}
\begin{itemize}
    \item Step 13 computes Scheme-5 reward: $r_t=\tanh\!\left(-\Delta\text{conflict}_{\text{rate}}\cdot\Delta f_{\text{pure,rate}}\cdot100\right)$
    \item Step 14 uses $r_t$ to update policy: better fitness → higher reward → reinforce current strategy
    \item This creates a feedback loop: RL learns which penalty adjustments lead to better optimization outcomes
\end{itemize}

\subsection{Step 15: Output Results}
Master process (rank 0) outputs comprehensive metrics:
\begin{equation}
\text{Output: } (\text{iter}, \text{eval}, h^{\text{penalty}}, f^{\text{pure}}, \text{penalty}, \text{improvement}, r, \text{conflict}, \alpha)
\end{equation}

\textbf{Detailed formulas for each output field:}

\begin{enumerate}
    \item \textbf{iter} (iteration number):
    \begin{equation}
    \text{iter} = t \quad \text{(current iteration counter, starts from 0)}
    \end{equation}
    
    \item \textbf{eval} (total function evaluations):
    \begin{equation}
    \text{eval} = \sum_{i=1}^N \text{eval}_i \quad \text{(sum of evaluations across all agents)}
    \end{equation}
    where $\text{eval}_i$ is the number of function evaluations performed by agent $i$ so far.
    
    \item \textbf{$h^{\text{penalty}}$} (average penalized fitness):
    \begin{equation}
    h^{\text{penalty}} = \frac{1}{N}\sum_{i=1}^N h_i(x_{i,\text{con}}) = \frac{1}{N}\sum_{i=1}^N \left[f_i(x_{i,\text{con}}) + \alpha_i(t) \cdot \text{conflict}_i(x_{i,\text{con}})\right]
    \end{equation}
    where $x_{i,\text{con}}$ is the consensus reference solution for agent $i$ (from Step 9).
    
    \item \textbf{$f^{\text{pure}}$} (average pure fitness):
    \begin{equation}
    f^{\text{pure}} = \frac{1}{N}\sum_{i=1}^N f_i(x_{i,\text{con}})
    \end{equation}
    This is the \textbf{target metric} to minimize (no penalty term).
    
    \item \textbf{penalty} (consensus violation measure):
    \begin{equation}
    \text{penalty} = h^{\text{penalty}} - f^{\text{pure}} = \frac{1}{N}\sum_{i=1}^N \left[\alpha_i(t) \cdot \text{conflict}_i(x_{i,\text{con}})\right]
    \end{equation}
    
    \textbf{Critical distinction: Penalty term vs. Penalty value}
    
    \begin{itemize}
        \item \textbf{Penalty term} (used in Step 1 optimization):
        \begin{equation}
        \text{Penalty term for agent } i \text{ at particle } x: \quad \alpha_i(t) \cdot \text{conflict}_i(x)
        \end{equation}
        This is applied to \textbf{each particle} during PSO optimization (Step 1):
        \begin{equation}
        h_i(x) = f_i(x) + \underbrace{\alpha_i(t) \cdot \text{conflict}_i(x)}_{\text{penalty term}}
        \end{equation}
        
        \item \textbf{Penalty value} (reported in Step 15 output):
        \begin{equation}
        \text{penalty} = \frac{1}{N}\sum_{i=1}^N \left[\alpha_i(t) \cdot \text{conflict}_i(x_{i,\text{con}})\right]
        \end{equation}
        This is the \textbf{average penalty term} across all agents, evaluated at the \textbf{consensus solution} $x_{i,\text{con}}$.
        
        \item \textbf{Key relationship:}
        \begin{itemize}
            \item \textbf{Penalty term} (per agent, per particle): $\alpha_i(t) \cdot \text{conflict}_i(x)$ - Used during optimization (Step 1)
            \item \textbf{Penalty value} (global average): $\text{penalty} = \frac{1}{N}\sum_{i=1}^N [\alpha_i(t) \cdot \text{conflict}_i(x_{i,\text{con}})]$ - Reported in output (Step 15)
            \item \textbf{Connection}: Penalty value = Average of all agents' penalty terms evaluated at consensus solution $x_{i,\text{con}}$
        \end{itemize}
        
        \item \textbf{Why they differ:}
        \begin{itemize}
            \item During optimization (Step 1), each agent evaluates penalty term for \textbf{many particles} ($M=300$ particles)
            \item In output (Step 15), we report penalty value for the \textbf{single consensus solution} $x_{i,\text{con}}$
            \item At convergence: Both penalty term and penalty value approach zero as $\text{conflict}_i \to 0$
        \end{itemize}
        
        \item         \textbf{Numerical example:}
        \begin{itemize}
            \item Agent 1: $\alpha_1 = 0.002$, $\text{conflict}_1(x_{1,\text{con}}) = 0.05$ $\Rightarrow$ penalty term $= 0.002 \times 0.05 = 0.0001$
            \item Agent 2: $\alpha_2 = 0.002$, $\text{conflict}_2(x_{2,\text{con}}) = 0.03$ $\Rightarrow$ penalty term $= 0.002 \times 0.03 = 0.00006$
            \item Penalty value (output) $= \frac{1}{2}(0.0001 + 0.00006) = 0.00008$
        \end{itemize}
    \end{itemize}
    
    \item \textbf{improvement} (relative fitness improvement):
    \begin{equation}
    \text{improvement}_t = \begin{cases}
    \frac{f_{t-1}^{\text{pure}} - f_t^{\text{pure}}}{|f_{t-1}^{\text{pure}}|} & \text{if } |f_{t-1}^{\text{pure}}| > \epsilon \quad \text{(relative improvement)} \\
    0 & \text{otherwise (first iteration or zero fitness)}
    \end{cases}
    \end{equation}
    where $\epsilon = 10^{-8}$ prevents division by zero. Positive improvement means fitness decreased (better solution).
    
    \item \textbf{$r$} (RL reward signal):
    \begin{equation}
    r_t = \tanh\!\left(-\Delta\text{conflict}_{\text{rate}} \cdot \Delta f_{\text{pure,rate}} \cdot 100\right)
    \end{equation}
    This is the same reward computed in Step 13, used for RL training in Step 14.
    
    \item \textbf{conflict} (average conflict measure):
    \begin{equation}
    \text{conflict} = \frac{1}{N}\sum_{i=1}^N \text{conflict}_i(x_{i,\text{best}})
    \end{equation}
    where $\text{conflict}_i(x_{i,\text{best}})$ is computed in Step 3:
    \begin{equation}
    \text{conflict}_i(x_{i,\text{best}}) = \sum_{d \in O_i} \frac{|x_{i,\text{best}}^d - x_{i,\text{con}}^d|}{x_{\max} - x_{\min}}
    \end{equation}
    Note: This uses $x_{i,\text{best}}$ (local best from Step 2), not $x_{i,\text{con}}$ (consensus reference).
    
    \textbf{Typical values:}
    \begin{itemize}
        \item Early iterations: $\text{conflict} \approx 50 \sim 150$ (high disagreement)
        \item Mid-stage: $\text{conflict} \approx 5 \sim 20$ (partial consensus)
        \item Late-stage: $\text{conflict} \to 0$ (convergence to consensus)
    \end{itemize}
    
    \item \textbf{$\alpha$} (current penalty coefficient):
    \begin{equation}
    \alpha = \frac{1}{N}\sum_{i=1}^N \alpha_i(t)
    \end{equation}
    In practice, if all agents use the same RL policy, $\alpha_i(t) = \alpha(t)$ for all $i$, so $\alpha = \alpha(t)$.
    The value $\alpha_i(t)$ is updated by RL in Step 14:
    \begin{equation}
    \alpha_i(t+1) = \text{clip}(\alpha_i(t) + a_t \times 0.01, 0.01, 1.0)
    \end{equation}
    where $a_t = \pi_\theta(s_t)$ is the action from the RL policy network.
\end{enumerate}

\textbf{Summary table of all output values:}

\begin{center}
\begin{tabular}{|l|l|l|}
\hline
\textbf{Field} & \textbf{Formula} & \textbf{Source Step} \\
\hline
iter & $t$ & Counter \\
\hline
eval & $\sum_{i=1}^N \text{eval}_i$ & Accumulated evaluations \\
\hline
$h^{\text{penalty}}$ & $\frac{1}{N}\sum_{i=1}^N h_i(x_{i,\text{con}})$ & Step 11 \\
\hline
$f^{\text{pure}}$ & $\frac{1}{N}\sum_{i=1}^N f_i(x_{i,\text{con}})$ & Step 11 \\
\hline
penalty & $\frac{1}{N}\sum_{i=1}^N [\alpha_i(t) \cdot \text{conflict}_i(x_{i,\text{con}})]$ & Step 11 (derived) \\
\hline
improvement & $\frac{f_{t-1}^{\text{pure}} - f_t^{\text{pure}}}{|f_{t-1}^{\text{pure}}|}$ & Step 12 \\
\hline
$r$ & $\tanh\!\left(-\Delta\text{conflict}_{\text{rate}}\cdot\Delta f_{\text{pure,rate}}\cdot100\right)$ & Step 13 \\
\hline
conflict & $\frac{1}{N}\sum_{i=1}^N \text{conflict}_i(x_{i,\text{best}})$ & Step 3 (averaged) \\
\hline
$\alpha$ & $\frac{1}{N}\sum_{i=1}^N \alpha_i(t)$ & Step 14 \\
\hline
\end{tabular}
\end{center}

\textbf{Important observations:}
\begin{itemize}
    \item \textbf{penalty} and \textbf{conflict} are different:
    \begin{itemize}
        \item \textbf{conflict}: Uses $x_{i,\text{best}}$ (local best), measures disagreement before consensus
        \item \textbf{penalty}: Uses $x_{i,\text{con}}$ (consensus reference), measures remaining violation after consensus
    \end{itemize}
    \item At convergence: Both $\text{conflict} \to 0$ and $\text{penalty} \to 0$, and $h^{\text{penalty}} = f^{\text{pure}}$
    \item The \textbf{penalty value} (output) is the average of penalty terms evaluated at consensus solution, not the penalty terms used during optimization
\end{itemize}

\subsection{Step 16: Convergence Check}
Termination condition:
\begin{equation}
\text{Stop if: } \text{eval\_count} \geq E_{\max}
\end{equation}

\subsection{No-Water-Down Guarantee (Theoretical)}
At convergence, consensus is achieved and penalized fitness equals pure fitness:
\begin{equation}
x_{i,\text{con}}^{*,d} = x_{j,\text{con}}^{*,d}, \forall d \in O_i \cap O_j \Rightarrow \text{conflict}_i = 0 \Rightarrow h_i(x_i^*) = f_i(x_i^*)
\end{equation}

\textbf{Proof sketch:}
\begin{align}
\text{At convergence: } & x_{i,\text{con}}^{*,d} = x_{j,\text{con}}^{*,d}, \forall d \in O_i \cap O_j \\
\Rightarrow & \text{conflict}_i(x_{i,\text{con}}^*) = \sum_{d \in O_i} \frac{|x_{i,\text{con}}^{*,d} - x_{i,\text{con}}^{*,d}|}{x_{\max} - x_{\min}} = 0 \\
\Rightarrow & h_i(x_i^*) = f_i(x_i^*) + \alpha_i(t) \cdot 0 = f_i(x_i^*) \\
\Rightarrow & h^{\text{penalty}} = \frac{1}{N}\sum_{i=1}^N h_i(x_i^*) = \frac{1}{N}\sum_{i=1}^N f_i(x_i^*) = f^{\text{pure}} \\
\Rightarrow & F(x^*) = \sum_{i=1}^N f_i(x_i^*) = N \cdot f^{\text{pure}}
\end{align}

\textbf{Practical implication:} The reported $f^{\text{pure}}$ is a valid evaluation of the global objective $F(x)$ divided by $N$, with no artificial inflation from penalty terms.

\end{document}

